\section{The Advantages You Cannot See}

\subsection{Asymmetries of Vision}

2 people attend the same conference on artificial intelligence. They hear the same talks, see the same demonstrations, and return to the same city. 1 starts an AI application for hospital billing. The other starts a verification tool for semiconductors. Neither could easily have invented the other's idea. The difference was not raw intelligence. It was what each already knew, which problems each recognized, whose complaints each had heard, and what each could plausibly build or sell.

This is the peculiar reality of entrepreneurship: opportunities are not waiting on a shelf, equally visible to anyone who looks. The world presents the same event to different people, and they notice different gaps. An experienced teacher sees the invisible scaffolding required to get students from watching a video to solving a new problem. A logistics manager sees inventory and routing mistakes everyone else takes for granted. A security engineer sees a vulnerability beneath a polished demo. The ability to notice a commercially consequential problem is itself an asset.

It~is tempting to imagine that launching an enterprise begins with an inspired idea. But the real beginning is more mundane. People accumulate knowledge, habits, contacts, and operational capabilities over years. A shift in technology or customer behavior makes some of these assets newly potent. 1 person recognizes an acute problem, gains direct access to its sufferers, and has enough ability to attempt a solution. The business emerges from that collision.

\subsection{The Diagnostic Question}

This chapter asks a question that sounds simple but turns out to be difficult: \textbf{What do you possess today that gives you a credible advantage in building some businesses rather than others?}

Not what you would like to be good at, nor what your résumé suggests, nor what you could theoretically master with unlimited time. What can you actually deploy (legally, reliably, in the coming weeks or months) to discover an unmet need, earn trust, create a result, find a buyer, and survive long enough to improve?

This is an exercise in economic diagnosis, not self-esteem. The output should be an empirical map of demonstrated assets, unknowns, constraints, plausible opportunities, and bounded experiments that can update the map. You should finish knowing what to test next, not merely feeling more confident.

The central insight is that \textbf{founder advantage is specific to a task and a market.} A brilliant mathematician may be an ineffective founder of a neighborhood bistro. A capable programmer who spent 5 years handling obscure insurance claims may be unusually well positioned to build useful insurance software. A person with no remarkable technical skill may possess uncommon buyer access, operational taste, or the ability to recruit the specialized people a venture needs.

What matters is the combination, and whether customers value what that combination can produce.
